% Generated draft: requires mathematical and source review.
\documentclass[11pt]{article}
\usepackage[margin=25mm]{geometry}
\usepackage{fontspec}
\setmainfont{Latin Modern Roman}
\usepackage{amsmath,amssymb,mathtools,unicode-math}
\setmathfont{Latin Modern Math}
\usepackage{enumitem,needspace,xurl,hyperref}
\setlength{\parindent}{0pt}
\setlength{\parskip}{5pt}
\setlength{\emergencystretch}{3em}
\newcommand{\N}{\mathbb{N}}
\newcommand{\Z}{\mathbb{Z}}
\newcommand{\Q}{\mathbb{Q}}
\newcommand{\R}{\mathbb{R}}
\newcommand{\C}{\mathbb{C}}
\newcommand{\F}{\mathbb{F}}
\newcommand{\A}{\mathbb{A}}
\newcommand{\PP}{\mathbb{P}}
\newcommand{\E}{\mathbb{E}}
\newcommand{\eps}{\varepsilon}
\newcommand{\dd}{\,\mathrm{d}}
\begin{document}
\section*{UnsolvedMath ID: 20003322}
\textbf{Source problem code:} catalogue entry
\subsection{Definitions and mathematical statement}
The phrase ``as above'' refers to section 5.6, recovered from companion permanent ID $20003321$. Let
\[
 Q(x_1,x_2)=ax_1^2+bx_1x_2+cx_2^2,\qquad a,b,c\in\Q,
\]
be nondegenerate, indefinite, and anisotropic over $\Q$, meaning $Q(x)=0$ for $x\in\Q^2$ only when $x=0$. Its discriminant $\Delta=b^2-4ac$ is positive and nonsquare in $\Q$, and its splitting field is $K=\Q(\sqrt\Delta)$. On $\mathbb T^2=\R^2/\Z^2$, define
\[
 m(Q,x)=\inf_{z\in\Z^2}|Q(x+z)|,\qquad
 M=m(Q)=\sup_{x\in\mathbb T^2}m(Q,x).
\]
Isolation of $M$ means that some $\delta>0$ satisfies
\[
 \{m(Q,x):x\in\mathbb T^2\}\cap(M-\delta,M)=\varnothing.
\]
After this isolation is established, define
\[
 m_2(Q)=\sup\{m(Q,x):x\in\mathbb T^2,\ m(Q,x)<M\}.
\]
Conjecture 36, attributed to E. Barnes and H. Swinnerton-Dyer and communicated by A. Pollington, states that $M\in\Q$ is isolated and that there exist torus classes with representatives $x,y\in K^2$ such that
\[
 m(Q,x)=M,\qquad m(Q,y)=m_2(Q).
\]
The supplied verification explicitly warns, citing Godwin, that $m_2(Q)$ need not itself be isolated. It also records the duplicated word ``the'' as a source typographical duplication, not OCR damage.

The attached aid reports a conditional finite certificate. First take $Q$ primitive integral of the above type and a hyperbolic integral automorphism $T\in\mathrm{SL}_2(\Z)$ with $Q(Tz)=Q(z)$. If for some $0<c<M$ the torus superlevel
\[
 X_c=\{x\in\mathbb T^2:m(Q,x)\ge c\}
\]
is finite, then $T$ permutes $X_c$. Every member is therefore periodic under $T$ and has a rational representative; its inhomogeneous minimum is a discrete rational value. The aid reports that this gives rational attainment and isolation of $M$. If the largest value below $M$ represented in $X_c$ exists and is strictly larger than $c$, it is $m_2(Q)$ and is attained at a rational representative. The positive gap between threshold and top is explicit here; finiteness of a set containing only points at a threshold equal to $M$ would not by itself certify isolation.

The reported normalization laws are
\[
 m(rQ,x)=|r|m(Q,x)\quad(r\in\Q\setminus\{0\}),\qquad
 m(Q\circ U,x)=m(Q,Ux)\quad(U\in\mathrm{GL}_2(\Z)),
\]
with the corresponding scaling and invariance of $m(Q)$ and, when defined, $m_2(Q)$. They transfer the certificate to rational forms. The aid gives the exact check
\[
 m(X^2-2Y^2)=\tfrac12,\qquad m(X^2-2Y^2,(0,\tfrac12))=\tfrac12.
\]
No finite superlevel certificate is supplied for every admissible form, so this conditional result does not establish the full conjecture.
\subsection{Short English statement}
For an anisotropic indefinite rational binary quadratic form, is its top inhomogeneous minimum rational and isolated, with both first and second minima attained over its splitting field? The supplied finite-superlevel certificate proves these conclusions only when its stated finite-set and threshold conditions hold; it does not require the second minimum to be isolated.
\subsection{Sources}
\begin{itemize}
\item E. Barnes and H. Swinnerton-Dyer [11], Conjecture 36, communicated by A. Pollington, American Institute of Mathematics, \emph{Emerging applications of measure rigidity}, section 5.6, \url{https://aimath.org/WWN/measrigid/measrigid.pdf}. The supplied verification mentions Godwin's warning about the second minimum; full bibliography entries are not supplied.
\item Companion local record, permanent ID $20003321$, source code \texttt{AIM-OTHER-0128}, canonical \texttt{aim-other-notes.json}, \texttt{notes[127]}, supplies the appended section 5.6 definitions and hypotheses, separate from its Question 35.
\item Local UnsolvedMath registry, permanent ID $20003322$, source code \texttt{catalogue entry}, canonical \texttt{aim-other-notes.json}, \texttt{notes[128]}, dated 14 August 2026, statement status corrected and verified according to supplied metadata. Research attempt 1: partial result, proved rigor, low novelty confidence; machine-generated and not peer reviewed.
\end{itemize}
\end{document}
